\section{Phương pháp}

\subsection{Research Approach}

Chúng tôi sử dụng \textbf{systematic literature review} kết hợp với \textbf{comparative analysis} để:
\begin{enumerate}
    \item Thu thập papers từ 2020--2026 trên 12 research domains
    \item Trích xuất evidence theo template chuẩn hóa
    \item So sánh quantitative results giữa các approaches
    \item Xác định research gaps và contradictory findings
    \item Đề xuất architecture dựa trên evidence tổng hợp
\end{enumerate}

\subsection{Search Strategy}
\begin{itemize}
    \item \textbf{Databases}: arXiv, ACL Anthology, NeurIPS, ICML, ICLR, AAAI, CHI, KDD
    \item \textbf{Keywords}: "AI character", "persona consistency", "memory architecture", "relationship agent", "emotional AI", "long-term interaction", "personality drift"
    \item \textbf{Inclusion criteria}: Papers 2020--2026, English/Vietnamese, peer-reviewed hoặc preprint có citation
    \item \textbf{Tiêu chí loại}: Lĩnh vực không liên quan, mô tả phương pháp không đủ
\end{itemize}

\subsection{Evidence Grading}

\begin{table}[H]
\centering
\caption{Cấp độ chất lượng bằng chứng}
\label{tab:evidence-quality}
\begin{tabular}{ll}
\toprule
\textbf{Cấp độ} & \textbf{Tiêu chí} \\
\midrule
Mạnh & Nhiều nghiên cứu khẳng định, N lớn, có ý nghĩa thống kê \\
Trung bình & Một nghiên cứu mạnh hoặc replication hạn chế \\
Yếu & Phát hiện sơ bộ, cỡ mẫu nhỏ, không kiểm định thống kê \\
Mâu thuẫn & Mâu thuẫn trực tiếp yêu cầu giải quyết \\
\bottomrule
\end{tabular}
\end{table}

% =============================================================================
% 5. LITERATURE REVIEW